\section{Antecedentes y Justificación}
\label{sec:antecedentes}

\subsection{Antecedentes}

Actualmente en la FICH cada carrera gestiona sus procesos de titulación e investigaciones de forma independiente, sin un sistema unificado. No existe una plataforma institucional ni externa que centralice el seguimiento de tesis y proyectos de grado. En contraste, otras facultades de la UAGRM como Postgrado, Escuela de Ingeniería, Ciencias Jurídicas, Políticas, Sociales y Relaciones Internacionales cuentan con sus propios sistemas, lo que evidencia una carencia en la FICH.

A nivel nacional, \textcite{perales2022repositorios} analizaron el desarrollo de repositorios y revistas científicas de acceso abierto en Bolivia, destacando los avances y rezagos en la comunicación científica universitaria. A nivel internacional, \textcite{leon2012implantacion} presentaron la implantación de un repositorio institucional en la Universidad de Sevilla, basado en tres ejes: recolección, preservación y divulgación. Ambos trabajos sirven como referencia para el diseño del sistema propuesto.

Si bien existen plataformas tecnológicas para la gestión académica, el presente proyecto se enfoca en el desarrollo de una solución a medida para la FICH, por lo que no se adopta ninguna solución existente.

\subsection{Justificación}

El presente proyecto se justifica por sus aportes en tres dimensiones. En el aspecto teórico, contribuye al fortalecimiento de los conceptos de gestión académica, trazabilidad documental y memoria institucional, aplicados al contexto del seguimiento de tesis y proyectos de grado. Desde el punto de vista metodológico, el enfoque de desarrollo basado en casos de uso y la modelización de flujos de trabajo académico constituyen una referencia replicable para proyectos similares en otras unidades de la UAGRM. En el ámbito práctico, el sistema reportará beneficios directos a la FICH al centralizar los procesos de titulación, mejorar la eficiencia operativa, reducir los tiempos de gestión, garantizar el control y respaldo documental, y consolidar la memoria institucional. Asimismo, impactará positivamente en los estudiantes al brindarles un seguimiento claro de su proyecto, en los docentes al facilitar una evaluación sistemática, y en las autoridades al proporcionar datos para la toma de decisiones.

\endinput
